\centerline{\bf \Huge Вспоминаем комбу I}

\begin{flushright}
        \scriptsize{– Дядь, скажи, который листочек?\\
                    – Первый, без теории.\\
                    – У тебя листочек-то комбинаторный чтоль?\\
                    – Так я же комбинаторщик.\\
                    – Да ну?\\
                    – Че, не веришь? Честное слово.\\
                    – Дядь, не в масть тебе такой листочек. Давай сдавай.\\
                    – А ты мне плюсик? \\
                    – Мена?\\
                    – Мена.\\
                    – Я тебе плюсик. Ты мне решение.\\
                    - Пошло}
\end{flushright}

\problem{1}
Назовем билет с номером от 000000 до 999999 отличным, если разность некоторых двух соседних цифр его номера равна 5. Найдите число отличных билетов.

\problem{2}
Сколькими способами можно выложить в ряд пять красных, пять синих и пять зелёных шаров так, чтобы никакие два синих шара не лежали рядом?

\problem{3}
Рассмотрим клетчатую таблицу $20 \times 25$. Посчитайте общее количество прямоугольников внутри неё. Например, внутри таблицы $2 \times 2$ ровно 9 прямоугольников: 4 прямоугольника $1 \times 1,4$ прямоугольника $1 \times 2,1$ прямоугольник $2 \times 2$.

\problem{4}
В Калмыкии 13 районов. Между некоторыми районами есть дорога. Известно, что всего этих дорог 67. Докажите, что Герман можно добраться из любого района в любой другой двигаясь по дорогам.

\problem{5}
Найдется ли 333 способа выбрать 3 различных числа от 1 до 18 таким образом, чтобы их произведение делилось на 9?

\problem{6}
Сколько существует 100-значных номеров, в которых хотя бы раз встречается тройка, а любые две цифры отличаются не больше чем на 1?